Vector search returns nearby embeddings from a large database. GPU
parallelism can accelerate distance calculations, but raw high-dimensional
vectors can exceed device memory. Compression can make larger indexes
resident, but search still has to read codes and select promising vectors.
JHQ is a hierarchical quantisation method that first compares compact
primary codes to form a shortlist, then uses residual codes for a more
detailed comparison \cite{jhq}. This thesis studies how to execute these
stages efficiently on a GPU while respecting their scoring and selection
rules, and measures the resulting quality--speed trade-off.

Graph-based GPU indexes store vectors and graph links; inverted-file
(IVF) indexes store the representations searched within selected database
partitions \cite{cagra,johnson_gpu,ivfrabitq}. If these structures exceed
device memory, the complete index cannot remain resident.
The largest benchmark, Stella, contains 17.78 million 1,024-dimensional
text embeddings. Its raw FP32 storage, available device memory and
JHQ-GPU resident allocation, respectively, illustrate the capacity
constraint (\cref{sec:7-7}):
\[
 S_{\rm raw}=67.81\,\mathrm{GiB}
 >S_{\rm GPU}\simeq31.8\,\mathrm{GiB}
 >S_{\rm resident}=20.44\,\mathrm{GiB}.
\]
This motivates a GPU-resident index that searches compressed codes
instead of keeping the full FP32 database on the device.
